\documentclass[a4paper,12pt]{article}
\usepackage[T2A]{fontenc}
\usepackage[utf8]{inputenc}
\usepackage[russian]{babel}
\usepackage{amsmath,amssymb,mathtools}
\usepackage{helvet}
\renewcommand{\familydefault}{\sfdefault}
\usepackage{geometry}
\geometry{left=20mm,right=20mm,top=21mm,bottom=21mm}
\usepackage{microtype}
\usepackage{booktabs,array,tabularx}
\usepackage{xcolor}
\usepackage{tikz}
\usepackage{hyperref}
\usepackage{fancyhdr}
\usepackage{parskip}
\usepackage{icomma}
\definecolor{accent}{HTML}{006E74}
\definecolor{deep}{HTML}{26343B}
\definecolor{err}{HTML}{A52C35}
\hypersetup{colorlinks=true,linkcolor=accent,pdfauthor={Лёвин Артём Александрович},pdftitle={Ревью домашней работы — Анастасия Павлова, 22 ноября 2026}}
\setlength{\parindent}{0pt}
\setlength{\parskip}{5pt plus 1pt}
\setlength{\headheight}{15pt}
\pagestyle{fancy}
\fancyhf{}
\fancyhead[L]{\small\color{deep} Вероятность и комбинаторика}
\fancyhead[R]{\small\color{deep} 22 ноября 2026 г.}
\fancyfoot[C]{\small\thepage}
\renewcommand{\headrulewidth}{0.3pt}
\renewcommand{\footrulewidth}{0pt}
\newcommand{\errlabel}{\textcolor{err}{\textbf{Ошибка:}}}
\newcommand{\keylabel}{\textcolor{accent}{\textbf{Ключевая идея:}}}
\newcommand{\outcomelabel}{\textcolor{accent}{\textbf{Итог:}}}
\newcommand{\reviewtask}[1]{\clearpage\phantomsection\label{task:#1}\section*{\textcolor{accent}{Задание №\,#1}}}
\begin{document}
\thispagestyle{empty}
\begin{center}
  \vspace*{34mm}
  {\Large\color{accent}\textbf{МАТЕМАТИКА}}\par
  \vspace{20mm}
  {\huge\bfseries Ревью домашней работы}\par
  \vspace{7mm}
  {\Large Вероятность и комбинаторика}\par
  \vspace{13mm}
  {\large Анастасия Павлова}\par
  \vspace{20mm}
  {\large 22 ноября 2026 г.}\par
  \vfill
  {\normalsize Лёвин Артём Александрович}\par
  {\small эксклюзивно для Анастасии Павловой}\par
  \vspace{12mm}
  \begin{tikzpicture}
    \draw[accent!65,line width=0.7pt] (-5.3,0) .. controls (-2.5,0.65) and (-1.4,-0.45) .. (0,0)
                .. controls (1.4,0.45) and (2.5,-0.65) .. (5.3,0);
  \end{tikzpicture}
  \vspace{3mm}
\end{center}

\clearpage
\section*{\textcolor{accent}{Сводная таблица замечаний}}
\small
\renewcommand{\arraystretch}{1.45}
\begin{tabularx}{\linewidth}{@{}p{1.0cm}p{3cm}X p{2.25cm}p{2.45cm}@{}}
\toprule
\textbf{№} & \textbf{Тема} & \textbf{Суть ошибки} & \textbf{Ваш ответ} & \textbf{Верный ответ}\\
\midrule
\hyperref[task:5]{\textbf{5}} & Ровно два орла & Не учтены три различных порядка двух орлов при трёх бросках. & $0{,}25$ & $\dfrac38=0{,}375$\\
\hyperref[task:6]{\textbf{6}} & Хотя бы одна решка & Количество бросков принято за число случаев; упущено дополнение к событию. & $0{,}75$ & $\dfrac{15}{16}=0{,}9375$\\
\hyperref[task:7]{\textbf{7}} & Формула Бернулли & Правильная дробь переведена в неверное десятичное число. & $0{,}27875$ & $\dfrac7{32}=0{,}21875$\\
\hyperref[task:10]{\textbf{10}} & Первый и последний участники & Учтён первый участник из России, но не учтено условие о последнем участнике из Китая. & $0{,}35$ & $\dfrac{63}{380}\approx0{,}1658$\\
\bottomrule
\end{tabularx}
\normalsize
\vspace{6mm}
В таблице приведены только задания, в которых обнаружены ошибки. По номерам заданий можно перейти к подробному разбору.

\clearpage
\thispagestyle{empty}
\vspace*{\fill}
\begin{center}
  {\LARGE\bfseries\color{accent} Разбор ошибок}
\end{center}
\vspace*{\fill}

\reviewtask{5}
\textbf{Условие.} Симметричную монету бросают трижды. Найдите вероятность того, что выпадет ровно два орла.

\textbf{Ваше решение.} На фотографии рассмотрена схема из трёх бросков; в качестве результата записано $\frac14=0{,}25$.

\textbf{Ваш ответ:} $0{,}25$.

\errlabel\ При подсчёте благоприятных исходов не учтены все порядки появления двух орлов. Один конкретный порядок выпадения орлов и решки составляет лишь один из возможных исходов.

\textbf{Где именно возникает ошибка.} До определения ответа верно выделены три последовательных броска монеты. Первый ошибочный вывод --- значение вероятности $\frac14$: оно не получается из равновозможных исходов трёх бросков при условии «ровно два орла».

\textbf{Почему это неверно.} Фраза «ровно два орла» допускает три порядка: орёл, орёл, решка; орёл, решка, орёл; решка, орёл, орёл. Все эти исходы необходимо считать благоприятными. Поскольку каждый бросок имеет два варианта, общее число последовательностей равно $2\cdot2\cdot2=8$. Эти восемь последовательностей равновозможны.

\keylabel\ В задачах про несколько бросков порядок важен. Для классической вероятности нужно посчитать \emph{все} благоприятные последовательности и разделить их число на количество всех последовательностей.

\textbf{Исправление от последнего правильного шага.} После определения трёх бросков перечислим все их исходы. Из восьми последовательностей ровно в трёх содержатся два орла.

\textbf{Полное правильное решение.} Обозначим орла буквой О, решку буквой Р. Возможные исходы:
\[
\underbrace{\text{ООО},\ \text{ООР},\ \text{ОРО},\ \text{ОРР},\
\text{РОО},\ \text{РОР},\ \text{РРО},\ \text{РРР}}_{8\ \text{равновозможных исходов}}.
\]
Ровно два орла появляются в последовательностях $\text{ООР},\ \text{ОРО},\ \text{РОО}$. Следовательно,
\[
P=\frac{3}{8}=0{,}375.
\]
\textbf{Проверка.} Можно выбрать место единственной решки тремя способами: первым, вторым или третьим. Каждый конкретный порядок имеет вероятность $\frac18$; поэтому суммарная вероятность $3\cdot\frac18=\frac38$.

\outcomelabel\ Вероятность выпадения ровно двух орлов равна $\boxed{\frac38=0{,}375}$.

\textbf{Задача на закрепление.} Симметричную монету бросают четыре раза. Найдите вероятность того, что выпадет ровно три орла.

\reviewtask{6}
\textbf{Условие.} Симметричную монету бросают четыре раза. Найдите вероятность того, что выпадет хотя бы одна решка.

\textbf{Ваше решение.} Записано $\frac34=0{,}75$.

\textbf{Ваш ответ:} $0{,}75$.

\errlabel\ Значение $\frac34$ не учитывает все исходы четырёх независимых бросков. Число бросков нельзя принимать за общее число равновозможных исходов.

\textbf{Где именно возникает ошибка.} В качестве знаменателя взято число $4$, хотя при четырёх бросках монеты всего $2^4=16$ равновозможных последовательностей. Числитель $3$ также не является числом последовательностей с хотя бы одной решкой.

\textbf{Почему это неверно.} При каждом броске возможны два результата. Результат первого броска сочетается с каждым результатом второго, затем третьего и четвёртого. Поэтому число исходов умножается: $2\cdot2\cdot2\cdot2=16$. Событие «хотя бы одна решка» включает последовательности с одной, двумя, тремя и четырьмя решками.

\keylabel\ Для события «хотя бы один раз произошло $A$» часто проще использовать противоположное событие «$A$ не произошло ни разу». Сумма их вероятностей равна единице.

\textbf{Исправление от последнего правильного шага.} Сохраняем условие о четырёх бросках, подсчитываем $16$ исходов и выделяем единственный исход без решек --- четыре орла подряд.

\textbf{Полное правильное решение.} Пусть $A$ --- событие «хотя бы одна решка». Тогда противоположное событие $\overline A$ --- «все четыре раза выпал орёл». Вероятность одного орла равна $\frac12$, броски независимы, поэтому
\[
P(\overline A)=\left(\frac12\right)^4=\frac1{16}.
\]
По правилу противоположных событий
\[
P(A)=1-P(\overline A)=1-\frac1{16}=\frac{15}{16}=0{,}9375.
\]
\textbf{Проверка.} Из всех $16$ последовательностей только одна не содержит решек. Значит, $15$ последовательностей благоприятны; классическая формула снова даёт $\frac{15}{16}$.

\outcomelabel\ Вероятность появления хотя бы одной решки равна $\boxed{\frac{15}{16}=0{,}9375}$.

\textbf{Задача на закрепление.} Симметричную монету бросают пять раз. Найдите вероятность того, что выпадет хотя бы один орёл.

\reviewtask{7}
\textbf{Условие.} Симметричную монету бросают восемь раз. По формуле Бернулли найдите вероятность выпадения ровно трёх орлов.

\textbf{Ваше решение.} Правильно записаны сочетание и вероятности:
\[
\frac{8!}{3!\,5!}\left(\frac12\right)^3\left(\frac12\right)^5
=56\cdot\frac1{2^8}=\frac{56}{256}.
\]
Однако далее на фотографии записано десятичное значение $0{,}27875$.

\textbf{Ваш ответ:} $0{,}27875$.

\errlabel\ Арифметическая ошибка при делении: дробь $\frac{56}{256}$ переведена в десятичную запись неверно.

\textbf{Где именно возникает ошибка.} Последний правильный шаг --- $\frac{56}{256}$. Первый неправильный шаг --- приравнивание этой дроби к $0{,}27875$.

\textbf{Почему это неверно.} Десятичную запись дроби получают делением числителя на знаменатель. Достаточно проверить найденное число обратным умножением: $0{,}27875\cdot256=71{,}36$, что не равно $56$. Следовательно, указанное десятичное значение не соответствует дроби.

\keylabel\ Даже при правильной формуле окончательное числовое значение нужно проверить. Перед делением удобно сократить дробь.

\textbf{Исправление от последнего правильного шага.} Сокращаем записанную Вами дробь на восемь:
\[
\frac{56}{256}=\frac{56:8}{256:8}=\frac7{32}.
\]
\textbf{Полное правильное решение.} Число бросков $n=8$, требуется $k=3$ орла, вероятность орла в одном броске $p=\frac12$. Формула Бернулли:
\[
P_n(k)=\binom nk p^k(1-p)^{n-k}.
\]
Подставим данные и последовательно вычислим:
\begin{align*}
P_8(3)
&=\binom83\left(\frac12\right)^3\left(\frac12\right)^5\\
&=\frac{8!}{3!\,5!}\cdot\frac1{2^8}
=56\cdot\frac1{256}\\
&=\frac7{32}=0{,}21875.
\end{align*}
\textbf{Проверка.} $0{,}21875\cdot32=7$, поэтому $\frac7{32}=0{,}21875$ точно. Полученная вероятность находится между нулём и единицей.

\outcomelabel\ Вероятность ровно трёх орлов равна $\boxed{\frac7{32}=0{,}21875}$.

\textbf{Задача на закрепление.} Симметричную монету бросают шесть раз. По формуле Бернулли найдите вероятность выпадения ровно двух орлов. Ответ запишите в виде несократимой дроби и десятичного числа.

\reviewtask{10}
\begingroup\small\setlength{\parskip}{2pt}
\textbf{Условие.} В соревновании участвуют 20 спортсменов: 7 из России, 9 из Китая и 4 из других стран. Порядок выступлений случаен. Найдите вероятность того, что первым выступит россиянин, а последним --- спортсмен из Китая.

\textbf{Ваше решение.} На фотографии приведён подсчёт
\[
\frac{7\cdot19!\cdot1}{20!}=\frac7{20}=0{,}35.
\]
\textbf{Ваш ответ:} $0{,}35$.

\errlabel\ Подсчитано событие «первым выступит россиянин», но дополнительное обязательное условие «последним выступит спортсмен из Китая» не включено в число благоприятных порядков.

\textbf{Где именно возникает ошибка.} Вы верно определили семь возможных российских участников на первом месте и $20!$ общих перестановок. Далее ошибочно использовано $19!$ способов расставить всех остальных. В $19!$ входят также порядки, где последним оказывается спортсмен из России или другой страны.

\textbf{Почему это неверно.} Нужны одновременно два требования, накладываемые на разные позиции. После выбора первого участника из России последний должен выбираться только из девяти участников из Китая. Остальные 18 спортсменов могут располагаться в любом порядке. При независимом случайном порядке места фиксируются последовательно; вероятность выбора китайского спортсмена на последнее место после фиксации первого россиянина составляет $\frac9{19}$.

\keylabel\ Когда в условии есть два обязательных требования, проверяйте, что оба учтены в подсчёте благоприятных исходов. Удобны два равноценных подхода: последовательные вероятности или подсчёт перестановок.

\textbf{Исправление от последнего правильного шага.} К семи способам выбрать первого добавляем девять способов выбрать последнего. Оставшихся спортсменов размещаем $18!$ способами:
\[
N_{\text{благ}}=7\cdot9\cdot18!.
\]
\textbf{Полное правильное решение (перестановки).} Всего имеется $20!$ равновозможных порядков выступлений. Для благоприятного порядка:
\begin{enumerate}
\item На первое место выбираем одного из семи россиян: $7$ способов.
\item На последнее место выбираем одного из девяти китайских спортсменов: $9$ способов.
\item Оставшихся 18 спортсменов размещаем между ними: $18!$ способов.
\end{enumerate}
Таким образом,
\begin{align*}
P&=\frac{7\cdot9\cdot18!}{20!}
=\frac{7\cdot9\cdot18!}{20\cdot19\cdot18!}\\
&=\frac{63}{380}\approx0{,}1658.
\end{align*}
\textbf{Проверка (последовательные вероятности).} Вероятность первого россиянина равна $\frac7{20}$. После закрепления этого участника остаётся 19 человек, из них 9 китайских спортсменов. Поэтому
\[
P=\frac7{20}\cdot\frac9{19}=\frac{63}{380}.
\]
Оба способа дают одно и то же значение. Число $0{,}35$ --- вероятность только первого условия; добавление второго условия должно уменьшать вероятность.

\outcomelabel\ Вероятность выполнить оба условия равна $\boxed{\frac{63}{380}\approx0{,}1658}$.

\textbf{Задача на закрепление.} В соревновании выступают 15 человек: 5 из команды А, 6 из команды Б и 4 из команды В. Порядок выступлений случаен. Найдите вероятность того, что первым выступит участник команды А, а последним --- участник команды Б.

\endgroup
\end{document}
